%%%%%%%%%%%%%%%%%%%%%%%%%%%%%%%%%%%%%%%%%%%%%%%%%%%%%%%%%%%%%%%%%%%%%%%%%%%%%%%
%% TKS FORMAL MATHEMATICAL MANUAL v6.0 — ACADEMIC RELEASE
%% Part 4 of 4: Compiler Spec, Theorems, Examples, Appendices
%%%%%%%%%%%%%%%%%%%%%%%%%%%%%%%%%%%%%%%%%%%%%%%%%%%%%%%%%%%%%%%%%%%%%%%%%%%%%%%

%% ============================================================================
%% PART IV: APPLICATIONS
%% ============================================================================
\part{Applications}

%% ============================================================================
%% CHAPTER 11: THE TKS LANGUAGE SPECIFICATION
%% ============================================================================
\chapter{The TKS Language: Syntax, Static Semantics, and Dynamic Semantics}
\label{ch:language-spec}

This chapter provides a complete language specification for TKS expressions, including grammar, parsing, type checking, and evaluation.

\section{Lexical Conventions}

\begin{definition}[Token Types]
\label{def:tokens}
\begin{verbatim}
TOKEN_WORLD     := 'A' | 'B' | 'C' | 'D'
TOKEN_MODE      := '1' | '2' | ... | '10'
TOKEN_NOETIC    := '^' followed by '0'..'9'
TOKEN_FOUND     := '_{' followed by '1'..'7' and optional 'a'..'d' and '}'
TOKEN_OP        := '⊕' | '⊖' | '⊕_T' | '⊖_T' | '⊗_T' | '⊘_T' | '∘' | '→'
TOKEN_FRACTAL_L := '⟨'
TOKEN_FRACTAL_R := '⟩'
TOKEN_LPAREN    := '('
TOKEN_RPAREN    := ')'
TOKEN_COLON     := ':'
\end{verbatim}
\end{definition}

\section{Complete BNF Grammar}

\begin{definition}[TKS Grammar (BNF)]
\label{def:bnf-grammar}
\begin{verbatim}
<program>       ::= <expression>+

<expression>    ::= <atom>
                  | <modified-expr>
                  | <compound-expr>
                  | <fractal-expr>
                  | '(' <expression> ')'

<atom>          ::= <element> | <acquisition> | <world>

<element>       ::= <world-sym> <mode>
<world-sym>     ::= 'A' | 'B' | 'C' | 'D'
<mode>          ::= '1' | '2' | '3' | '4' | '5' | '6' | '7' | '8' | '9' | '10'

<acquisition>   ::= 'A0' | <acq-type> <foundation-num>
<acq-type>      ::= 'D' | 'W' | 'P'
<foundation-num>::= '1' | '2' | '3' | '4' | '5' | '6' | '7'

<world>         ::= 'A' | 'B' | 'C' | 'D'

<modified-expr> ::= <expression> <noetic>? <foundation>?

<noetic>        ::= '^' <noetic-num>
<noetic-num>    ::= '0' | '1' | '2' | '3' | '4' | '5' | '6' | '7' | '8' | '9'

<foundation>    ::= '_{' <foundation-num> <sub-world>? '}'
<sub-world>     ::= 'a' | 'b' | 'c' | 'd'

<compound-expr> ::= <expression> <binary-op> <expression>

<binary-op>     ::= '⊕' | '⊖' | '⊕_T' | '⊖_T' | '⊗_T' | '⊘_T' | '∘' | '→'

<fractal-expr>  ::= '⟨' <noetic-list> '⟩' '(' <expression> ')'
<noetic-list>   ::= <noetic-num> (':' <noetic-num>)*
\end{verbatim}
\end{definition}

\section{Abstract Syntax Tree}

\begin{definition}[AST Node Types]
\label{def:ast-nodes}
\begin{verbatim}
AST_Node :=
  | ElementNode(world: World, mode: Mode)
  | NoeticNode(k: Int, child: AST_Node)
  | FoundationNode(m: Int, w: SubWorld?, child: AST_Node)
  | BinaryOpNode(op: Operator, left: AST_Node, right: AST_Node)
  | FractalNode(chain: List[Int], child: AST_Node)
  | WorldAssocNode(world: World, child: AST_Node)
\end{verbatim}
\end{definition}

\section{Static Semantics (Type Checking)}

\begin{definition}[Type Checking Algorithm]
\label{def:typecheck}
\begin{verbatim}
FUNCTION TypeCheck(ast: AST_Node, ctx: Context) → Type | Error:
  MATCH ast:
    CASE ElementNode(world, mode):
      IF mode < 1 OR mode > 10: RETURN TypeError("Invalid mode")
      RETURN Element(world, mode)

    CASE NoeticNode(k, child):
      child_type = TypeCheck(child, ctx)
      IF child_type IS Error: RETURN child_type
      IF k < 0 OR k > 9: RETURN NoeticError("Invalid Noetic index")
      RETURN Expr

    CASE FoundationNode(m, w, child):
      child_type = TypeCheck(child, ctx)
      IF child_type IS Error: RETURN child_type
      IF m < 1 OR m > 7: RETURN FoundationError("Invalid Foundation")
      IF w NOT IN {a,b,c,d,null}: RETURN FoundationError("Invalid sub-world")
      RETURN Expr

    CASE BinaryOpNode(op, left, right):
      left_type = TypeCheck(left, ctx)
      right_type = TypeCheck(right, ctx)
      IF left_type IS Error: RETURN left_type
      IF right_type IS Error: RETURN right_type
      RETURN CheckBinaryOp(op, left_type, right_type)

    CASE FractalNode(chain, child):
      IF length(chain) = 0: RETURN FractalError("Empty fractal")
      FOR k IN chain:
        IF k < 0 OR k > 9: RETURN NoeticError("Invalid Noetic in fractal")
      child_type = TypeCheck(child, ctx)
      IF child_type IS Error: RETURN child_type
      RETURN Expr
\end{verbatim}
\end{definition}

\section{Dynamic Semantics (Evaluation)}

\begin{definition}[Evaluation Order]
\label{def:eval-order}
The mandatory evaluation sequence is:
\begin{enumerate}
  \item Parse input to AST
  \item Type-check AST
  \item Apply Foundation filters (innermost first)
  \item Apply Noetic operators (innermost first)
  \item Evaluate Tootra arithmetic (by precedence)
  \item Apply World associations
  \item Expand Fractals
  \item Check RPM prerequisites
  \item Return final value or error
\end{enumerate}
\end{definition}

\begin{definition}[Big-Step Operational Semantics]
\label{def:bigstep}
We define the evaluation relation $\Downarrow$:

\textbf{Element Evaluation:}
\[
\frac{}{Wn \Downarrow e_{W,n}} \quad [\textsc{E-Elem}]
\]

\textbf{Noetic Evaluation:}
\[
\frac{E \Downarrow v}{\noe{k}(E) \Downarrow \mathfrak{n}_k(v)} \quad [\textsc{E-Noetic}]
\]

\textbf{Foundation Evaluation:}
\[
\frac{E \Downarrow v}{E_{mw} \Downarrow \varphi_{m,w}(v)} \quad [\textsc{E-Found}]
\]

\textbf{Addition Evaluation:}
\[
\frac{E_1 \Downarrow v_1 \quad E_2 \Downarrow v_2}{E_1 \Tadd E_2 \Downarrow v_1 \uplus v_2} \quad [\textsc{E-Add}]
\]

\textbf{Fractal Evaluation:}
\[
\frac{E \Downarrow v}{\langle k_1 : \cdots : k_n \rangle(E) \Downarrow \mathfrak{n}_{k_1}(\cdots\mathfrak{n}_{k_n}(v)\cdots)} \quad [\textsc{E-Fractal}]
\]
\end{definition}

\section{Normal Forms}

\begin{definition}[TKS Normal Form]
\label{def:normal-form}
An expression $E$ is in \textbf{normal form} iff:
\begin{enumerate}
  \item All Noetics are applied (no pending ${}^k$)
  \item All Foundations are resolved (no pending ${}_{mw}$)
  \item All Fractals are expanded
  \item All binary operations are evaluated
  \item World associations are finalized
  \item Expression is a single Element or composite Idea
\end{enumerate}
\end{definition}

\begin{definition}[Normalization Rules]
\label{def:norm-rules}
\begin{align}
(E^k)^\ell &\to E^{k\ell} && \text{(Stack Noetics)} \\
E^0 &\to E && \text{(Identity elimination)} \\
(E_{m_1 w_1})_{m_2 w_2} &\to E_{m_2 w_2} && \text{(Outer Foundation dominates)} \\
\langle k \rangle(E) &\to E^k && \text{(Single-element fractal)} \\
E^{23} &\to E^0 \approx E && \text{(Duality cancellation)} \\
E^{56} &\to E^0 \approx E && \text{(FM cancellation)} \\
E^{89} &\to E^0 \approx E && \text{(AB cancellation)}
\end{align}
\end{definition}

%% ============================================================================
%% CHAPTER 12: NOETIC FRACTAL CALCULUS
%% ============================================================================
\chapter{Noetic Fractal Calculus}
\label{ch:fractal-calculus}

This chapter provides a complete treatment of Noetic fractals, their algebra, and their convergence properties.

\section{Fractal Notation and Semantics}

\begin{definition}[Noetic Fractal]
\label{def:noetic-fractal}
A \textbf{Noetic Fractal} is a nested sequence of Noetic operators:
\[
\langle k_1 : k_2 : \cdots : k_n \rangle(E) = \noe{k_1}(\noe{k_2}(\cdots\noe{k_n}(E)\cdots))
\]
\end{definition}

\begin{definition}[Single-Nesting Fractal]
\label{def:single-nest}
The notation $X.Y$ denotes ``$Y$ within $X$'':
\[
X.Y = \langle X : Y \rangle = \noe{X} \circ \noe{Y}
\]
\end{definition}

\section{Complete Single-Nesting Fractal Table}

\begin{longtable}{lll}
\caption{All 100 Single-Nesting Fractals} \\
\toprule
\textbf{Fractal} & \textbf{Composition} & \textbf{Interpretation} \\
\midrule
\endfirsthead
\toprule
\textbf{Fractal} & \textbf{Composition} & \textbf{Interpretation} \\
\midrule
\endhead
\multicolumn{3}{c}{\textbf{Mind Fractals ($1.X$ and $X.1$)}} \\
\midrule
$1.0$ & $\noe{1} \circ \noe{0}$ & Awareness of pure potential \\
$1.1$ & $\noe{1} \circ \noe{1}$ & Meta-awareness \\
$1.2$ & $\noe{1} \circ \noe{2}$ & Awareness of attraction \\
$1.3$ & $\noe{1} \circ \noe{3}$ & Awareness of repulsion \\
$1.4$ & $\noe{1} \circ \noe{4}$ & Awareness of vibration (MVR core) \\
$1.5$ & $\noe{1} \circ \noe{5}$ & Awareness of reception \\
$1.6$ & $\noe{1} \circ \noe{6}$ & Awareness of projection \\
$1.7$ & $\noe{1} \circ \noe{7}$ & Awareness of rhythm (MVR core) \\
$1.8$ & $\noe{1} \circ \noe{8}$ & Awareness of elevation \\
$1.9$ & $\noe{1} \circ \noe{9}$ & Awareness of grounding \\
\midrule
$0.1$ & $\noe{0} \circ \noe{1}$ & Potential awareness \\
$2.1$ & $\noe{2} \circ \noe{1}$ & Attraction to awareness \\
$3.1$ & $\noe{3} \circ \noe{1}$ & Repulsion from awareness \\
$4.1$ & $\noe{4} \circ \noe{1}$ & Vibrating awareness \\
$5.1$ & $\noe{5} \circ \noe{1}$ & Receiving awareness \\
$6.1$ & $\noe{6} \circ \noe{1}$ & Projecting awareness \\
$7.1$ & $\noe{7} \circ \noe{1}$ & Rhythmic awareness \\
$8.1$ & $\noe{8} \circ \noe{1}$ & Elevated awareness \\
$9.1$ & $\noe{9} \circ \noe{1}$ & Grounded awareness \\
\midrule
\multicolumn{3}{c}{\textbf{Polarity Fractals}} \\
\midrule
$2.2$ & $\noe{2}^2$ & Deep attraction \\
$2.3$ & $\noe{2} \circ \noe{3}$ & Attraction to repulsion \\
$3.2$ & $\noe{3} \circ \noe{2}$ & Repulsion of attraction \\
$3.3$ & $\noe{3}^2$ & Deep repulsion \\
\midrule
\multicolumn{3}{c}{\textbf{Structure Fractals}} \\
\midrule
$5.5$ & $\noe{5}^2$ & Deep reception \\
$5.6$ & $\noe{5} \circ \noe{6}$ & Receiving projection \\
$6.5$ & $\noe{6} \circ \noe{5}$ & Projecting reception \\
$6.6$ & $\noe{6}^2$ & Deep projection \\
\midrule
\multicolumn{3}{c}{\textbf{Dynamic Fractals}} \\
\midrule
$4.4$ & $\noe{4}^2$ & Resonance \\
$4.7$ & $\noe{4} \circ \noe{7}$ & Vibrating rhythm \\
$7.4$ & $\noe{7} \circ \noe{4}$ & Rhythmic vibration \\
$7.7$ & $\noe{7}^2$ & Meta-rhythm \\
\midrule
\multicolumn{3}{c}{\textbf{Causation Fractals}} \\
\midrule
$8.8$ & $\noe{8}^2$ & Deep elevation \\
$8.9$ & $\noe{8} \circ \noe{9}$ & Elevation of grounding \\
$9.8$ & $\noe{9} \circ \noe{8}$ & Grounding of elevation \\
$9.9$ & $\noe{9}^2$ & Deep grounding \\
\bottomrule
\end{longtable}

\section{Fractal Algebra}

\begin{definition}[Fractal Composition]
\label{def:fractal-comp}
\[
\langle A : B \rangle \circ \langle C : D \rangle = \langle A : B : C : D \rangle
\]
\end{definition}

\begin{definition}[Fractal Simplification]
\label{def:fractal-simp}
When dual pairs appear adjacent:
\begin{align}
\langle \cdots : 2 : 3 : \cdots \rangle &\approx \langle \cdots : 0 : \cdots \rangle \\
\langle \cdots : 5 : 6 : \cdots \rangle &\approx \langle \cdots : 0 : \cdots \rangle \\
\langle \cdots : 8 : 9 : \cdots \rangle &\approx \langle \cdots : 0 : \cdots \rangle
\end{align}
\end{definition}

\begin{theorem}[Polarity Oscillation Collapse]
\label{thm:polarity-collapse}
\[
\langle 2 : 3 : 2 : 3 : 2 : 3 \rangle \approx \langle 0 \rangle = \noe{0}
\]
Alternating polarities cancel to neutral.
\end{theorem}

\begin{proof}
By repeated application of duality neutralization:
\[
\langle 2 : 3 \rangle^3 \approx (\noe{0})^3 = \noe{0}
\]
\end{proof}

\section{Convergence and Stability}

\begin{definition}[Stabilizing Fractal]
\label{def:stab-fractal}
A fractal $F$ is \textbf{stabilizing} if repeated application converges:
\[
\lim_{n \to \infty} F^n(X) = X^*
\]
for some fixed point $X^*$.
\end{definition}

\begin{definition}[Destabilizing Fractal]
\label{def:destab-fractal}
A fractal $F$ is \textbf{destabilizing} if:
\[
\|F^n(X)\| \to \infty \text{ as } n \to \infty
\]
\end{definition}

\begin{theorem}[MVR Stability]
\label{thm:mvr-stab}
The MVR fractal $\langle 1 : 4 : 7 \rangle$ is stabilizing for compatible inputs.
\end{theorem}

%% ============================================================================
%% CHAPTER 13: THEOREMS AND PROOFS
%% ============================================================================
\chapter{Theorems and Proofs}
\label{ch:theorems}

This chapter collects and proves the major theorems of TKS v6.0.

\section{Fundamental Structural Theorems}

\begin{theorem}[Noetica is a Category]
\label{thm:noetica-cat-formal}
The structure $\Noetica$ with objects $\Ob(\Noetica)$ and morphisms $\Hom(\Noetica)$ satisfies all category axioms.
\end{theorem}

\begin{proof}
See Theorem~\ref{thm:noetica-category} in Chapter~\ref{ch:category-noetica}.
\end{proof}

\begin{theorem}[ACBE is a Functor]
\label{thm:acbe-func-formal}
$\mathrm{ACBE} : \Noetica_A \to \Noetica_D$ is a valid functor preserving identity and composition.
\end{theorem}

\begin{proof}
See Theorem~\ref{thm:acbe-functor} in Chapter~\ref{ch:acbe-functor}.
\end{proof}

\begin{theorem}[RPM is a Monad]
\label{thm:rpm-monad-formal}
The structure $(\RPM, \eta, \mu)$ forms a monad on $\Acq$ satisfying the three monad laws.
\end{theorem}

\begin{proof}
See Theorem~\ref{thm:rpm-monad-laws} in Chapter~\ref{ch:rpm-monad}.
\end{proof}

\section{Termination and Decidability}

\begin{theorem}[RPM Termination]
\label{thm:rpm-term-formal}
The RPM diagnostic algorithm terminates for any input in $O(22)$ steps.
\end{theorem}

\begin{proof}
See Theorem~\ref{thm:rpm-term} in Chapter~\ref{ch:rpm-monad}.
\end{proof}

\begin{theorem}[Type Checking Decidability]
\label{thm:typecheck-decide}
Type checking for TKS expressions is decidable.
\end{theorem}

\begin{proof}
The type system has finitely many types, and the inference rules are syntax-directed. Each rule can be checked in constant time. The total complexity is $O(n)$ for an expression of size $n$.
\end{proof}

\section{Algebraic Theorems}

\begin{theorem}[Noetic Monoid]
\label{thm:noetic-monoid-formal}
$(\Noetics, \circ, \noe{0})$ forms a monoid.
\end{theorem}

\begin{proof}
See Theorem~\ref{thm:monoid} in Chapter~\ref{ch:noetic-algebra}.
\end{proof}

\begin{theorem}[Tootra Semiring]
\label{thm:tootra-semiring-formal}
$(\Ideas, \Tadd, \Tmul, \varnothing, \mathbf{1})$ forms a commutative semiring under simplified operations.
\end{theorem}

\begin{proof}
See Theorem~\ref{thm:semiring} in Chapter~\ref{ch:tootra-arithmetic}.
\end{proof}

\begin{theorem}[ACBE Preserves Causal Order]
\label{thm:acbe-order}
The ACBE transformation preserves the world ordering $A \succ B \succ C \succ D$.
\end{theorem}

\begin{proof}
ACBE factors as $F_D \circ F_C \circ F_B$, where each $F_W$ increases the world ordinal by 1. For $X \in \Noetica_A$ with $\ord(X) = 0$:
\[
\ord(\mathrm{ACBE}(X)) = \ord(F_D(F_C(F_B(X)))) = 3 = \ord(D)
\]
The direction $A \to B \to C \to D$ is preserved.
\end{proof}

\begin{theorem}[Stable Eigenmodes Exist]
\label{thm:eigenmodes}
For each Noetic $\noe{k}$, there exists at least one stable eigenstate with $\lambda = 1$.
\end{theorem}

\begin{proof}
For each $\noe{k}$ where $k \in \{1,\ldots,9\}$, the Element $Xk$ (for any World $X$) satisfies $\noe{k}(Xk) = Xk$. For $\noe{0}$, any Element is an eigenstate. Each Noetic has 4 stable eigenstates (one per World).
\end{proof}

%% ============================================================================
%% CHAPTER 14: WORKED EXAMPLES
%% ============================================================================
\chapter{Worked Examples}
\label{ch:examples}

This chapter provides fully worked examples demonstrating TKS concepts in action.

\section{Healing Chain Examples}

\begin{example}[Emotional Wound Healing]
\label{ex:emotional-healing}
\textbf{Expression:}
\[
(C3 \ominus C)^1 \circ (A2 \oplus C)^1_{3c}
\]

\textbf{Step-by-step evaluation:}

\textit{Step 1:} Parse the expression.
\begin{itemize}
  \item Left operand: $(C3 \ominus C)^1$
  \item Right operand: $(A2 \oplus C)^1_{3c}$
  \item Operation: Composition $\circ$
\end{itemize}

\textit{Step 2:} Evaluate left operand.
\begin{itemize}
  \item $C3$ = Emotional Negative (pain)
  \item $C3 \ominus C$ = Disassociate from Emotional world
  \item $(C3 \ominus C)^1$ = Apply Mind (conscious awareness)
  \item Result: Conscious removal of emotional pain
\end{itemize}

\textit{Step 3:} Evaluate right operand.
\begin{itemize}
  \item $A2$ = Spiritual Positive (virtue)
  \item $A2 \oplus C$ = Associate with Emotional world
  \item $(A2 \oplus C)^1$ = Apply Mind
  \item $(\ldots)_{3c}$ = Filter to Emotional Life foundation
  \item Result: Conscious infusion of virtue for emotional vitality
\end{itemize}

\textit{Step 4:} Apply composition.
\begin{itemize}
  \item Clear shadow first, then fill with light
  \item Result: Healed emotional field
\end{itemize}

\textbf{Final result:}
\[
\mathrm{EmotionalHealing} = \mathrm{ClearedField} \Tadd \mathrm{SpiritualVirtue}
\]
\end{example}

\begin{example}[Physical Health Restoration]
\label{ex:health}
\textbf{Expression:}
\[
\langle 1 : 4 : 7 \rangle(D2^2_{3d})
\]

\textbf{Evaluation:}

\textit{Step 1:} Evaluate base.
\[
D2^2_{3d} = \noe{2}(\varphi_{3,d}(D2)) = \text{Attract physical health for bodily vitality}
\]

\textit{Step 2:} Expand fractal.
\[
\langle 1 : 4 : 7 \rangle(\ldots) = \noe{1}(\noe{4}(\noe{7}(\ldots)))
\]

\textit{Step 3:} Apply from inside out.
\begin{itemize}
  \item $\noe{7}$: Establish rhythm (repetition)
  \item $\noe{4}$: Charge with vibration (intensity)
  \item $\noe{1}$: Apply mind (awareness)
\end{itemize}

\textbf{Result:}
\[
\mathrm{MVR}(D2^2_{3d}) = \text{Rhythmic, vibratory, conscious health pattern}
\]
\end{example}

\section{ACBE Manifestation Examples}

\begin{example}[Career Manifestation]
\label{ex:career}
\textbf{Expression:}
\[
\mathrm{ACBE}(A8^1_{6d})
\]

\textbf{Full cascade:}
\begin{center}
\begin{tabular}{llll}
\toprule
\textbf{Stage} & \textbf{Element} & \textbf{World} & \textbf{Interpretation} \\
\midrule
$A8^1_{6d}$ & $A8$ & Spiritual & Esoteric career truth \\
$\downarrow F_B$ & $B8$ & Mental & Profound career insight \\
$\downarrow F_C$ & $C8$ & Emotional & Passionate alignment \\
$\downarrow F_D$ & $D8$ & Physical & High quality action \\
\bottomrule
\end{tabular}
\end{center}

\textbf{Result:} Career manifested in Material foundation with high quality.
\end{example}

\section{RPM Diagnostic Examples}

\begin{example}[Wealth Failure Diagnosis]
\label{ex:wealth-failure}
\textbf{Scenario:} Person wants wealth but consistently fails.

\textbf{RPM Chain for $F_6$ (Material):}
\[
A0 \to D_6 \to W_6 \to P_6 \to \mathrm{Outcome}_6
\]

\textbf{Diagnostic:}
\begin{itemize}
  \item $A0$: $\checkmark$ Pure desire exists
  \item $D_6$: $\checkmark$ Wants material resources
  \item $W_6$: $\times$ \textbf{FAILED} --- Believes ``money is evil''
  \item $P_6$: $\triangle$ Blocked (depends on $W_6$)
\end{itemize}

\textbf{Failure Origin:} $W_6$ (Wisdom-Material)

\textbf{Correction:} Repair mental models about wealth before attempting action.
\end{example}

\section{Type Checking Examples}

\begin{example}[Well-Typed Expression]
\label{ex:welltyped}
\textbf{Expression:} $A8^1_{2b}$

\textbf{Derivation:}
\[
\frac{
  \frac{A : \mathtt{World} \quad 8 : \mathtt{Mode}}{A8 : \mathtt{Element}} [\textsc{Elem}]
  \quad 1 \in \{0,\ldots,9\}
}{A8^1 : \mathtt{Expr}} [\textsc{Noetic}]
\]
\[
\frac{A8^1 : \mathtt{Expr} \quad 2 \in \{1,\ldots,7\} \quad b \in \{a,b,c,d\}}{A8^1_{2b} : \mathtt{Expr}} [\textsc{Found}]
\]

\textbf{Result:} Well-typed as $\mathtt{Expr}$.
\end{example}

\begin{example}[Ill-Typed Expression]
\label{ex:illtyped}
\textbf{Expression:} $A11$

\textbf{Analysis:}
\[
\frac{A : \mathtt{World} \quad 11 : ?}{???}
\]

Since $11 \notin \{1,\ldots,10\}$, mode is out of range.

\textbf{Result:} $\mathtt{TypeError}$: ``Invalid mode 11''
\end{example}

\section{Fractal Transformation Examples}

\begin{example}[Deep Meditation Fractal]
\label{ex:meditation}
\textbf{Expression:}
\[
\langle 1:1:1:4:4:4:7:7:7 \rangle(A1_{1a})
\]

\textbf{Interpretation:} Triple $\mathrm{MVR}^3$ for extremely stable spiritual awareness.

\textbf{Evaluation:}
\[
= (\noe{1})^3 \circ (\noe{4})^3 \circ (\noe{7})^3 (A1_{1a})
\]

Each component is applied three times for deep installation.
\end{example}

\begin{example}[Polarity Balancing]
\label{ex:polarity}
\textbf{Expression:}
\[
\langle 2:3:2:3:2:3 \rangle(C4_{1c})
\]

\textbf{Simplification:}
\[
\langle 2:3 \rangle^3 \approx (\noe{0})^3 = \noe{0}
\]

\textbf{Result:} Returns to balanced neutral state.
\end{example}

%% ============================================================================
%% CHAPTER 15: DISCUSSION AND FUTURE WORK
%% ============================================================================
\chapter{Discussion and Future Work}
\label{ch:discussion}

\section{Summary of Contributions}

TKS v6.0 provides:
\begin{enumerate}
  \item A rigorous category-theoretic foundation ($\Noetica$, ACBE functor, RPM monad)
  \item Complete denotational semantics for all TKS operations
  \item Formal type theory with inference rules and safety properties
  \item Full algebraic specification of Noetic and Tootra operations
  \item Complete BNF grammar and evaluation semantics
  \item Expanded fractal calculus with convergence analysis
\end{enumerate}

\section{Relationship to Other Fields}

TKS connects to:
\begin{itemize}
  \item \textbf{Category Theory}: TKS is fundamentally categorical
  \item \textbf{Type Theory}: Full type system with inference
  \item \textbf{Programming Language Semantics}: Denotational and operational semantics
  \item \textbf{Formal Epistemology}: Prerequisite modeling
  \item \textbf{Cognitive Science}: Consciousness modeling
\end{itemize}

\section{Future Directions}

Potential extensions include:
\begin{enumerate}
  \item Automated theorem proving for TKS
  \item Implementation of TKS interpreter/compiler
  \item Connection to topos theory
  \item Probabilistic extensions
  \item Quantum TKS
\end{enumerate}

%% ============================================================================
%% APPENDICES
%% ============================================================================
\appendix
\part{Appendices}

\chapter{Symbol Reference}
\label{app:symbols}

\begin{longtable}{lll}
\toprule
\textbf{Symbol} & \textbf{Name} & \textbf{Meaning} \\
\midrule
$\Worlds$ & World set & $\{A, B, C, D\}$ \\
$\Ideas$ & Idea space & All Ideas \\
$\Elements$ & Element space & 40 Elements \\
$\Noetics$ & Noetic set & $\{\noe{0}, \ldots, \noe{9}\}$ \\
$\Foundations$ & Foundation set & $\{F_1, \ldots, F_7\}$ \\
$\Acquisitions$ & Acquisition set & 22 Acquisitions \\
$\Noetica$ & Category Noetica & TKS category \\
$\Acq$ & Category Acq & Acquisition category \\
$\RPM$ & RPM Monad & Prerequisite monad \\
$\Tadd$ & Tootra-Addition & Idea union \\
$\Tsub$ & Tootra-Subtraction & Idea removal \\
$\Tmul$ & Tootra-Multiplication & Idea interaction \\
$\Tdiv$ & Tootra-Division & Idea decomposition \\
$\sem{\cdot}$ & Semantic brackets & Denotation \\
$\types$ & Type judgment & ``has type'' \\
\bottomrule
\end{longtable}

\chapter{Quick Reference Cards}
\label{app:quickref}

\section{Element Grid}

\[
\begin{array}{c|cccccccccc}
 & 1 & 2 & 3 & 4 & 5 & 6 & 7 & 8 & 9 & 10 \\
\hline
A & A1 & A2 & A3 & A4 & A5 & A6 & A7 & A8 & A9 & A10 \\
B & B1 & B2 & B3 & B4 & B5 & B6 & B7 & B8 & B9 & B10 \\
C & C1 & C2 & C3 & C4 & C5 & C6 & C7 & C8 & C9 & C10 \\
D & D1 & D2 & D3 & D4 & D5 & D6 & D7 & D8 & D9 & D10 \\
\end{array}
\]

\section{Acquisition Matrix}

\[
\begin{array}{c|ccccccc}
 & F_1 & F_2 & F_3 & F_4 & F_5 & F_6 & F_7 \\
\hline
D & D_1 & D_2 & D_3 & D_4 & D_5 & D_6 & D_7 \\
W & W_1 & W_2 & W_3 & W_4 & W_5 & W_6 & W_7 \\
P & P_1 & P_2 & P_3 & P_4 & P_5 & P_6 & P_7 \\
\end{array}
\]
Root: $A0$ (Pure Desire)

\section{Mirror Principle}

Noetic pairs summing to 9:
\begin{align*}
\noe{1} + \noe{8} &= 9 && \text{Mind} \leftrightarrow \text{Above} \\
\noe{2} + \noe{7} &= 9 && \text{Positive} \leftrightarrow \text{Rhythm} \\
\noe{3} + \noe{6} &= 9 && \text{Negative} \leftrightarrow \text{Male} \\
\noe{4} + \noe{5} &= 9 && \text{Vibration} \leftrightarrow \text{Female}
\end{align*}

\section{Dual Pairs (Pseudo-Inverses)}

\begin{align*}
\noe{2} &\leftrightarrow \noe{3} && \text{Positive} \leftrightarrow \text{Negative} \\
\noe{5} &\leftrightarrow \noe{6} && \text{Female} \leftrightarrow \text{Male} \\
\noe{8} &\leftrightarrow \noe{9} && \text{Above} \leftrightarrow \text{Below}
\end{align*}

\chapter{Glossary}
\label{app:glossary}

\begin{description}
\item[ACBE] Above-Cause-Below-Effect; the manifestation functor from spiritual to physical.
\item[Acquisition] One of 22 prerequisite states ($A0$ + $7 \times D$ + $7 \times W$ + $7 \times P$).
\item[Element] One of 40 fundamental units formed by World $\times$ Mode.
\item[Foundation] One of 7 life domains (Unity, Wisdom, Life, Companionship, Power, Material, Lust).
\item[Fractal] Nested Noetic composition, written $\langle X : Y : Z \rangle$.
\item[Idea] Any content that can be processed by Mind; the fundamental operand.
\item[Mind] The fundamental operator that processes Ideas.
\item[MVR] Mind-Vibration-Rhythm; the core installation protocol.
\item[Noetic] One of 10 transformation operators ($\noe{0}$ through $\noe{9}$).
\item[Noetica] The category formed by TKS objects and morphisms.
\item[RPM] Recursive Prerequisite Model; the diagnostic monad.
\item[Sub-Foundation] Foundation $\times$ World; one of 28 contextual filters.
\item[Tootra Arithmetic] The four operations $\Tadd$, $\Tsub$, $\Tmul$, $\Tdiv$.
\item[World] One of 4 metaphysical planes (Spiritual, Mental, Emotional, Physical).
\end{description}

\chapter{Version History}
\label{app:version}

\begin{center}
\begin{tabular}{lll}
\toprule
\textbf{Version} & \textbf{Date} & \textbf{Changes} \\
\midrule
1.0 & --- & Initial formulation \\
2.0 & --- & Added Foundations and Sub-Foundations \\
3.0 & --- & Added RPM and ACBE \\
3.2.1 & --- & Formalized Category Noetica \\
5.0 & 2024 & Complete unification; Category Theory; Type Theory \\
\textbf{6.0} & \textbf{2024} & \textbf{Academic Release: Denotational semantics,} \\
 & & \textbf{strengthened proofs, complete type theory,} \\
 & & \textbf{full compiler specification} \\
\bottomrule
\end{tabular}
\end{center}

%% ============================================================================
%% BACK MATTER
%% ============================================================================
\backmatter

\chapter*{Colophon}
\addcontentsline{toc}{chapter}{Colophon}

\begin{center}
\textbf{TKS FORMAL MATHEMATICAL MANUAL v6.0}

\textit{Academic Release}

The Complete Rigorous Formalization of the Tootra Kabbalistic System

\vspace{2em}

This document represents the complete formal specification of TKS, suitable for academic study, formal verification, and computational implementation.

All canonical definitions preserved from v5.0.\\
All major upgrades completed as specified.

\vspace{2em}

\textbf{END OF DOCUMENT}
\end{center}

\end{document}
